%=============================================================================!
% 2.E  EXERCISES                                                              !
%=============================================================================!
\unnumbered{Exercises}
\label{sec:ne-exercises}

The exercises run from questions of understanding, through calculations,
to short derivations marked `show that'. Full solutions are in
\hyperref[sec:ne-solutions]{`Solutions to the exercises'} at the end of
the book, and Notebook~02 checks every numerical answer.

\begin{exercise}[a lift]
\label{ex:ne-lift}
A lift rises at a steady \SI{2}{\metre\per\second}. What is the net force
on a passenger standing in it? Is a frame of reference fixed to the lift
inertial? What changes while the lift is speeding up?
\end{exercise}

\begin{exercise}[a cyclist]
\label{ex:ne-frames}
A cyclist riding at \SI{6}{\metre\per\second} throws a ball forwards at
\SI{4}{\metre\per\second} relative to the bicycle. How fast does the ball
move relative to the ground? How fast if it is thrown backwards at the
same speed?
\end{exercise}

\begin{exercise}[two forces]
\label{ex:ne-net}
A body of mass \SI{3}{\kilogram} is pulled by a force of \SI{4}{\newton}
due east and a force of \SI{3}{\newton} due north. Find the size and
direction of the net force and of the acceleration.
\end{exercise}

\begin{exercise}[measuring mass]
\label{ex:ne-mass}
The same stretched spring gives a \SI{2}{\kilogram} trolley an
acceleration of \SI{0.6}{\metre\per\second\squared} and a second trolley
an acceleration of \SI{1.5}{\metre\per\second\squared}. Find the mass of
the second trolley and the force of the spring.
\end{exercise}

\begin{exercise}[a horse and cart]
\label{ex:ne-horse}
A horse pulls a cart, and by the third law the cart pulls back on the
horse with an equal force. Explain why the two can still speed up
together. (Draw a free-body diagram for each, and include the push of the
ground on the horse's hooves.)
\end{exercise}

\begin{exercise}[a rope]
\label{ex:ne-rope}
A \SI{2}{\kilogram} box is lifted by a vertical rope with an upward
acceleration of \SI{1.5}{\metre\per\second\squared}. Draw its free-body
diagram and find the tension in the rope.
\end{exercise}

\begin{exercise}[on the Moon]
\label{ex:ne-moon}
On the Moon, $g = \SI{1.62}{\metre\per\second\squared}$. Find the mass
and the weight of a \SI{70}{\kilogram} astronaut there, and compare the
weight with that on the Earth.
\end{exercise}

\begin{exercise}[a state]
\label{ex:ne-state}
At $t = 0$ a ball is \SI{3}{\metre} above the ground and moving downwards
at \SI{2}{\metre\per\second}. Ignoring air resistance, when does it reach
the ground? (Setting the height to zero gives a quadratic equation in $t$,
whose positive root is the answer; compare \cref{fig:ne-state}.)
\end{exercise}

\begin{exercise}[shorter steps]
\label{ex:ne-euler}
Show that one step of length $\tau$ from the start of the throw in
\cref{sec:ne-equations} overshoots the true height by exactly
$\tfrac12 g\tau^2$. Show that two steps of length $\tau/2$, covering the
same time, overshoot by $\tfrac14 g\tau^2$, half as much.
\end{exercise}

\begin{exercise}[faster than terminal]
\label{ex:ne-above}
The bead of \cref{sec:ne-drag}, with $\gamma = \SI{5}{\per\second}$, is
fired downwards into the liquid at \SI{4}{\metre\per\second}. Find its
velocity after \SI{0.2}{\second}.
\end{exercise}

\begin{exercise}[the approach]
\label{ex:ne-approach}
For a bead released from rest with $\gamma = \SI{5}{\per\second}$, find
the time at which its velocity reaches \SI{99}{\percent} of the terminal
velocity.
\end{exercise}

\begin{exercise}[checking a solution]
\label{ex:ne-drag-check}
Show by differentiating \cref{eq:ne-drag-x} that $\dd x/\dd t$ is the
velocity $v_\infty(1 - \mathrm{e}^{-\gamma t})$, and that $\dd^2 x/\dd
t^2 = g - \gamma\,\dd x/\dd t$, which is \cref{eq:ne-drag}.
\end{exercise}

\begin{exercise}[a stiff spring]
\label{ex:ne-spring}
A \SI{0.2}{\kilogram} block on a spring of stiffness
\SI{80}{\newton\per\metre} is pulled \SI{3}{\centi\metre} from its
resting place and released. Find the angular frequency, the period, the
largest speed and the largest acceleration.
\end{exercise}

\begin{exercise}[a kick]
\label{ex:ne-kick}
A block with $\omega = \SI{10}{\radian\per\second}$ rests at $x = 0$ and
is struck so that it starts with $v_0 = \SI{0.5}{\metre\per\second}$.
Find $A$ and $\phi$, and show that the motion is $x = A\sin\omega t$.
\end{exercise}

\begin{exercise}[a hanging spring]
\label{ex:ne-hanging}
A block of mass $m$ hangs at rest from a vertical spring of stiffness
$k$. Show that the spring is stretched by $mg/k$. Measuring $x$
downwards from this resting place, show that the block obeys $\dd^2
x/\dd t^2 = -(k/m)\,x$, so that gravity shifts the centre of the
oscillation but leaves its period unchanged.
\end{exercise}

\begin{exercise}[coupling]
\label{ex:ne-coupling}
On a level, smooth track, a \SI{2}{\kilogram} cart moving at
\SI{3}{\metre\per\second} runs into a stationary \SI{1}{\kilogram} cart
and the two couple together. Using momentum, find their common velocity
afterwards.
\end{exercise}

\begin{exercise}[walking on a boat]
\label{ex:ne-boat}
A person of \SI{70}{\kilogram} stands at one end of a \SI{140}{\kilogram}
boat at rest on still water, and walks \SI{3}{\metre} along the boat
towards the other end. Neglecting the drag of the water, use the centre of mass to find
how far the boat moves, and in which direction.
\end{exercise}

\begin{exercise}[an oxygen atom]
\label{ex:ne-oxygen}
An oxygen atom (\SI{15.999}{\amu}) at rest feels a constant force of
\SI{0.5}{\electronvolt\per\angstrom}. Find its acceleration in
\si{\angstrom\per\femto\second\squared} and its velocity after
\SI{5}{\femto\second}, in \si{\angstrom\per\femto\second} and in
\si{\metre\per\second}.
\end{exercise}

\begin{exercise}[removing the drift]
\label{ex:ne-drift}
Show that subtracting $\vb{P}/M$ from the velocity of every atom leaves a
total momentum of zero, and that it changes no acceleration.
\end{exercise}
